\subsection{Constancy in the original analytic family}
The preceding results now establish the numerical hypothesis needed by
simultaneous resolution.  The passage through algebraic families has served
only to compare the scalar invariant; we return afterwards to the original
holomorphic family.

\begin{theorem}[Constancy of the logarithmic invariant]
\label{thm:beta-constant}
Let $F(x,t)$ be a holomorphic family of isolated hypersurface singularities
in $\mathbb C^3$ with $f_t(0)=0$ and constant positive Milnor number.
If the central normal surface germ $X_0=\{f_0=0\}$ is not log canonical,
then $\beta(X_t,0)=\beta(X_0,0)$ for every sufficiently small $t$.
\end{theorem}

\begin{proof}
By Proposition~\ref{prop:scalar-curve-reduction}, it suffices to compare
a polynomial family over a smooth algebraic curve at a marked non-log-canonical
fiber with its nearby general fibers.  Log canonicity is constant by
Lemma~\ref{lem:lc-persistence}, so all germs in the relevant connected
curve are non-log-canonical.  Apply Lemma~\ref{lem:compactification} to
obtain a projective hypersurface family with exactly these singularities
and smooth complement of the section.  Its punctured-fiber Euler
characteristic and its canonical square are constant.

Construct the model of Proposition~\ref{prop:log-model}, making its finite
smooth curve base change.  Write
\[
 Y_0=S_0+\sum_a S_a,\qquad L=K_Y+H,
 \qquad S^\circ=S\setminus\operatorname{Supp}B_S.
\]
Lemma~\ref{lem:surface-boundary} and
Proposition~\ref{prop:full-exceptional-boundary} verify all the
adjunction and full-boundary requirements.  In particular,
\[
 S_0^\circ\simeq X_0\setminus\Gamma,
 \qquad Y_c\setminus H\simeq X_c\setminus\Gamma\quad(c\ne0).
\]
The first of these open surfaces is smooth.  On every other component,
any isolated interior singularity is smoothed by the family, and
Lemma~\ref{lem:interior-correction} gives a nonnegative correction
$\chi(F_x)-1$ to its nearby Euler value.

Apply Proposition~\ref{prop:nearby-euler-comparison}.  Constancy of the
punctured-fiber Euler characteristic and the two displayed isomorphisms
cancel the contribution of $S_0^\circ$, giving
\begin{equation}\label{eq:zero-euler-sum}
 0=\sum_a\chi_c(S_a^\circ)+\delta,
 \qquad
 \delta=\sum_a\sum_{x\in\operatorname{Sing}(S_a)\cap S_a^\circ}
                  (\chi(F_x)-1)\ge0.
\end{equation}
Each $S_a$ is projective and its adjoint divisor is nef, so
Lemma~\ref{lem:surface-bmy} gives
\[
 0\le(K_{S_a}+B_{S_a})^2\le3\chi_c(S_a^\circ).
\]
Every term on the right side of \eqref{eq:zero-euler-sum} is therefore
nonnegative.  Every extra logarithmic square is zero.

All hypotheses of Proposition~\ref{prop:beta-from-log-model} have now
been verified.  It gives equality of the logarithmic invariant on the
marked and nearby general fibers.  The finite base change used in the
model repeats the individual surface germs, so this is equality for
the original algebraic family as well.  Proposition~\ref{prop:scalar-curve-reduction}
then yields the assertion for the given holomorphic family.
\end{proof}
